% !TEX program = xelatex
% 공통수학2 · Ⅱ. 집합과 명제 · 1. 집합 · 중단원 마무리 (25차시)
\documentclass[11pt,a4paper]{article}
\usepackage{kotex}
\usepackage{iftex}
\ifPDFTeX\else
  \setmainhangulfont{Noto Serif CJK KR}
  \setsanshangulfont{Noto Sans CJK KR}
\fi
\usepackage[margin=2.1cm]{geometry}
\usepackage{parskip}
\usepackage{enumitem}
\usepackage{array}
\usepackage{tabularx}
\usepackage{booktabs}
\usepackage{xcolor}
\usepackage{amssymb}
\usepackage{tikz}
\usepackage[most]{tcolorbox}
\usepackage{fancyhdr}
\usepackage{titlesec}

\definecolor{goldc}{HTML}{B07C22}
\definecolor{goldstrong}{HTML}{8A5F16}
\definecolor{indigoc}{HTML}{2B3B57}
\definecolor{indigosoft}{HTML}{6E82A6}
\definecolor{linec}{HTML}{D6DACB}
\definecolor{surface2}{HTML}{E4E8DC}
\definecolor{notebg}{HTML}{FBF3E3}
\definecolor{inksoft}{HTML}{52604F}

\titleformat{\section}{\normalfont\sffamily\bfseries\large\color{indigoc}}{\thesection}{0.6em}{}
\titlespacing{\section}{0pt}{16pt}{8pt}
\newtcolorbox{ideabox}{enhanced, breakable, colback=white, colframe=goldc, boxrule=0.5pt, leftrule=3pt, arc=2pt, left=10pt, right=10pt, top=8pt, bottom=8pt}
\newtcolorbox{calloutbox}{enhanced, breakable, colback=white, colframe=indigosoft, boxrule=0.5pt, leftrule=3pt, arc=2pt, left=10pt, right=10pt, top=7pt, bottom=7pt}
\newtcolorbox{answerbox}{enhanced, breakable, colback=white, colframe=linec, boxrule=0.5pt, arc=2pt, left=10pt, right=10pt, top=7pt, bottom=7pt}
\newcommand{\lessonbanner}[3]{%
  \begin{tcolorbox}[enhanced, colback=indigoc, colframe=indigoc, boxrule=0pt, arc=0pt,
    left=14pt, right=14pt, top=14pt, bottom=10pt, borderline south={2.4pt}{0pt}{goldc}, fontupper=\color{white}]
    {\sffamily\footnotesize\color{white!85!black} #1}\\[4pt]{\sffamily\bfseries\Large #2}\\[3pt]{\sffamily\small\color{white!88!black} #3}
  \end{tcolorbox}}
\pagestyle{fancy}\fancyhf{}\renewcommand{\headrulewidth}{0pt}
\fancyfoot[L]{\footnotesize\color{inksoft} 공통수학2 · 2026학년도 2학기 · 지평선고등학교}
\fancyfoot[R]{\footnotesize\color{inksoft} 교사용 지도안 -- 배포 금지}

\begin{document}
\lessonbanner{공통수학2 · Ⅱ. 집합과 명제}{1. 집합 --- 중단원 마무리 (25차시)}{교사용 지도안 (2026학년도 2학기)}
\vspace{10pt}

\noindent\begin{tabularx}{\linewidth}{@{}>{\bfseries\color{indigoc}}p{2.6cm}X@{}}
\toprule
대상 & 고1 · 2개 반 · 반당 13명 \\
차시 & 50분 (도입 5 · 전개 35 · 정리 10) \\
목적 & 18$\sim$24차시(집합, 포함관계, 교집합·합집합, 여집합·차집합)를 종합 정리한다. \\
\bottomrule
\end{tabularx}

\section{개념 지도 총정리}
\begin{ideabox}
집합 중단원 전체의 흐름: \textbf{기준을 정해 모임을 정의(집합) $\to$ 두 집합 비교(포함 관계) $\to$ 두 집합을 조합(교집합·합집합) $\to$ 기준에서 제외(여집합·차집합)}. 벤 다이어그램은 이 모든 개념을 시각적으로 통합하는 도구였다.
\end{ideabox}
\begin{tabularx}{\linewidth}{@{}>{\bfseries\small}p{4.6cm}X@{}}
\toprule
\textbf{개념} & \textbf{핵심 공식} \\
\midrule
부분집합 & $A\subset B$: $A$의 모든 원소가 $B$에 속함 \\
교집합·합집합 & $A\cap B=\{x\mid x\in A$ 그리고 $x\in B\}$, $A\cup B=\{x\mid x\in A$ 또는 $x\in B\}$ \\
원소의 개수 & $n(A\cup B)=n(A)+n(B)-n(A\cap B)$ \\
여집합·차집합 & $A^C=\{x\mid x\in U$ 그리고 $x\notin A\}$, $A-B=A\cap B^C$ \\
드모르간의 법칙 & $(A\cup B)^C=A^C\cap B^C$, $(A\cap B)^C=A^C\cup B^C$ \\
\bottomrule
\end{tabularx}

\section{수업 흐름}
\begin{tabularx}{\linewidth}{@{}>{\centering\arraybackslash}p{1.6cm}X>{\raggedright\arraybackslash\small}p{3.1cm}@{}}
\toprule
\textbf{단계} & \textbf{활동 \& 발문} & \textbf{ATL / VTR} \\
\midrule
\textcolor{goldstrong}{\textbf{도입}}\newline\footnotesize 5분 &
\textbf{마음공부 (2분)} --- 중단원 전체 유무념 대조.\newline
\textbf{개념 지도 확인 (3분)} --- 위 표 빈칸형 복습. &
ATL: 자기관리 \\
\addlinespace
\textcolor{goldstrong}{\textbf{전개}}\newline\footnotesize 35분 &
\textbf{기본 (15분)} --- 마무리 문제 01$\sim$05(집합 판별, 원소나열법, 포함관계, 원소개수, 종합 집합연산) 개별 풀이 후 교차 채점.\newline
\textbf{표준 (10분)} --- 06, 07(서로소인 집합 나열, 조건을 만족하는 집합 X의 개수)을 모둠 협력으로 해결.\newline
\textbf{도전 (10분)} --- 서술형 11과 발전 09, 10, 12(자연수 배수 집합, 드모르간의 법칙 활용, 최댓값·최솟값) 중 학급 수준에 맞게 선택. &
ATT: 촉진자·평가자\newline ATL: 협업·자기점검 \\
\addlinespace
\textcolor{goldstrong}{\textbf{정리}}\newline\footnotesize 10분 &
\textbf{자기 평가 (5분)} --- 성취수준 체크리스트.\newline
\textbf{성찰 (5분)} --- 감사 일기 + 다음 중단원(명제) 예고. &
마음공부 · 자기평가 \\
\bottomrule
\end{tabularx}

\begin{calloutbox}
\textbf{지도상의 유의점} --- 이 차시는 진단·정리의 시간이다. 오답이 나온 문제가 18$\sim$24차시 중 어느 차시 개념인지 학생 스스로 표시하여 보충 지점을 시각화한다. 문제 09, 12는 드모르간의 법칙 활용이 핵심이므로 시간이 부족하면 우선순위로 다룬다.
\end{calloutbox}

\section{형성평가 답안 (지도서 원문 01$\sim$12번)}
\begin{tabularx}{\linewidth}{@{}>{\bfseries\color{goldstrong}}p{1.4cm}X@{}}
\toprule
01 & ⑴ $\times$(기준 불명확) ⑵ $\bigcirc$($1,2,3$) ⑶ $\times$(기준 불명확) \\
02 & ⑴ $\{1,3,7,21\}$ ⑵ $\{1,6\}$ \\
03 & ⑴ $A\subset B$ ⑵ $B\subset A$ \\
04 & $n(A\cap B)=18+23-30=11$ \\
05 & $U=\{1,\dots,10\}$, $A=\{1,5,9,10\}$, $B=\{1,2,5,10\}$: ⑴ $A\cap B=\{1,5,10\}$ ⑵ $A\cup B=\{1,2,5,9,10\}$ ⑶ $A^C=\{2,3,4,6,7,8\}$ ⑷ $A-B=\{9\}$ \\
06 & $A=\{1,2,4\}$와 서로소인 $U=\{1,\dots,5\}$의 부분집합: $\varnothing,\{3\},\{5\},\{3,5\}$ \\
07 & $X\subset(B-A)=\{12,16,20\}$의 부분집합 $\to$ 개수 $8$ \\
08 & $A\cap B=B$일 조건에서 $x=-1,0,4$ \\
09 & 분배법칙 활용 $\to$ $n(A_{10}\cap(A_3\cup A_8))=5$ \\
10 & 드모르간의 법칙 $\to$ $n(A^C\cap B^C)=20$ \\
11 & 최댓값 $21$, 최솟값 $9$ \\
12 & $n(A^C\cup B^C)=99-3=96$ \\
\bottomrule
\end{tabularx}

\section{다음 중단원}
\begin{calloutbox}
\textbf{2. 명제 (26차시$\sim$).} 집합에서 다룬 `참·거짓을 판별할 수 있는 명확한 기준'이, 이번엔 `문장의 참·거짓'을 다루는 명제로 확장된다. 사실 조건과 진리집합은 이미 배운 집합 개념과 정확히 대응된다.
\end{calloutbox}
\end{document}
